\documentclass[12pt]{article}
\usepackage[utf8]{inputenc}
\usepackage[spanish]{babel}
\usepackage{amsmath}
\usepackage{amsfonts}
\usepackage{amssymb}
\usepackage{geometry}
\geometry{letterpaper, margin=2.5cm}

\begin{document}

\begin{center}
    \Large \textbf{Evaluación de Examen 1: Unidad I} \\
    \large Física para Ciencias de la Salud (005-1713) \\
    \vspace{0.5cm}
    \textbf{Estudiante:} Alanyermus Pérez \quad \textbf{C.I.:} 33.171.611
    \vspace{0.3cm} \\
    \large \textbf{Calificación Final: 8/10 puntos}
\end{center}

\vspace{0.5cm}

\section*{Análisis de Resultados}

\subsection*{1. Ángulo plano (Conversión a radianes)}
\textbf{Procedimiento y resultado:} El estudiante utiliza una regla de tres basándose en la equivalencia $90^\circ = \frac{\pi}{2}$ rad. Aunque la notación en los pasos intermedios es algo informal, la lógica matemática aplicada es correcta. Llega a la respuesta de $0,41\pi$ rad, la cual es una aproximación válida del valor exacto $\frac{5\pi}{12} \approx 0,416\pi$. \\
\textbf{Veredicto:} \textbf{Correcto} (2/2 pts).

\subsection*{2. Sistemas de coordenadas (Longitud del hueso)}
\textbf{Procedimiento y resultado:} Excelente resolución. El estudiante dibuja el gráfico ubicando los puntos correctamente, identifica las distancias de los catetos del triángulo rectángulo formado en el plano ($a=6$, $b=8$) y aplica el Teorema de Pitágoras de manera impecable para obtener $L = 10$ cm. \\
\textbf{Veredicto:} \textbf{Correcto} (2/2 pts).

\subsection*{3. Vectores (Fuerza resultante)}
\textbf{Procedimiento y resultado:} Plantea correctamente que la fuerza resultante es la suma de los vectores $\vec{F}_1$ y $\vec{F}_2$. Agrupa bien las componentes $\hat{i}$ y $\hat{j}$. Sin embargo, comete un \textbf{error de signo} en la suma de la componente $\hat{j}$: en lugar de sumar $(25 + 35)\hat{j}$, escribe $(25 - 35)\hat{j}$. Esto lo lleva al resultado erróneo de $\vec{F} = 30\hat{i} - 10\hat{j}$ N (cuando el vector correcto es $30\hat{i} + 60\hat{j}$ N). \\
\textbf{Veredicto:} \textbf{Parcialmente correcto} (Planteamiento inicial válido, error de signo al operar) (1/2 pts).

\subsection*{4. Potencias de diez (Notación científica y prefijo)}
\textbf{Procedimiento y resultado:} Expresa la medida $0,0000125$ m empleando la potencia $12,5 \times 10^{-6}$. Aunque la notación científica estándar suele requerir un coeficiente entre 1 y 10 (es decir, $1,25 \times 10^{-5}$), el uso de la potencia a la $-6$ (notación de ingeniería) facilitaba la conversión. No obstante, el estudiante \textbf{omitió la segunda parte de la instrucción}: no indicó el prefijo adecuado del Sistema Internacional, que correspondía a la palabra ``micrómetros'' ($\mu$m). \\
\textbf{Veredicto:} \textbf{Incompleto / Parcialmente correcto} (1/2 pts).

\subsection*{5. Sistemas de unidades (Conversión de densidad)}
\textbf{Procedimiento y resultado:} El estudiante tachó un primer intento con factores de conversión formales, pero debajo realizó el cálculo aritmético correcto de forma procedimental: dividió el valor numérico entre $1000$ (asociado al paso de g a kg) y luego multiplicó por $1\,000\,000$ (asociado al paso de $\text{cm}^3$ a $\text{m}^3$ en el denominador). El resultado final numérico de $1030 \text{ kg/m}^3$ es exacto. \\
\textbf{Veredicto:} \textbf{Correcto} (2/2 pts).

\vspace{0.5cm}
\section*{Comentarios Generales y Sugerencias}
El estudiante demuestra un buen dominio general de los conceptos básicos evaluados (factores de conversión, Teorema de Pitágoras y el concepto de suma vectorial). 
\begin{itemize}
    \item Se sugiere recomendar al estudiante prestar mayor atención a la ley de los signos durante la suma algebraica (como ocurrió en la Pregunta 3).
    \item Se recomienda recordarle leer con mayor detenimiento los enunciados completos para asegurar que responde a todas las partes solicitadas (como ocurrió al omitir el prefijo del S.I. en la Pregunta 4).
\end{itemize}

\end{document}